% !TeX root = Javier_TARAZONA_CV_fr.tex
% !TeX program = pdflatex
%-------------------------
% Resume in Latex
% Author : Javier Tarazona
% Offre : BNP Paribas - Stage - AI Research Intern H/F (ITG Research, ref S_ITGIDF27_0088) - bwelcome.hr.bnpparibas (Avature)
%------------------------

\documentclass[a4paper,10pt]{article}
\DeclareUnicodeCharacter{00A0}{~}
\DeclareUnicodeCharacter{2013}{\textendash}
\DeclareUnicodeCharacter{2014}{\textemdash}
\DeclareUnicodeCharacter{201C}{``}
\DeclareUnicodeCharacter{201D}{''}
\DeclareUnicodeCharacter{2022}{\textbullet}
\DeclareUnicodeCharacter{25E6}{\textopenbullet}
\usepackage[T1]{fontenc}      % Codificación de salida (acentos y símbolos)
\usepackage[utf8]{inputenc}   % Lee texto en UTF-8
\usepackage[empty]{fullpage}
\usepackage[pdftex]{hyperref}
\usepackage[usenames,dvipsnames]{color}
\usepackage{enumitem}         % Control de listas
\usepackage{lmodern}
\usepackage{titlesec}
\usepackage{xcolor}

% --- Viñetas limpias y seguras ---
\setlist[itemize]{label=\textbullet, labelsep=0.5em, leftmargin=*}

% ===== CORRECTIFS ATS =====
\input{glyphtounicode}        % chaque glyphe -> vrai texte Unicode (ligatures decomposees)
\pdfgentounicode=1
\usepackage[none]{hyphenat}   % point 5 : plus aucun mot-cle coupe
\sloppy                       %           evite les debordements de marge
\usepackage{textcomp}         % point 6 : apostrophe ASCII au lieu de U+2019
\catcode`\'=\active
\def'{\textquotesingle}
\hypersetup{                  % point 8 : metadonnees indexables
  pdftitle={Javier Andres TARAZONA JIMENEZ - CV - AI Research Intern - Stage de fin d'études en Intelligence Artificielle},
  pdfauthor={Javier Andres TARAZONA JIMENEZ},
  pdfsubject={AI Research Intern - stage de fin d'études Bac+5 de 6 mois à partir d'avril 2027 - recherche en Intelligence Artificielle, Machine Learning, Deep Learning, Large Language Models (LLM)},
  pdfkeywords={AI Research Intern, Intelligence Artificielle, recherche, Machine Learning, Deep Learning,
               Large Language Models (LLM), architectures Transformer, NLP, prompting, fine-tuning,
               apprentissage supervisé, apprentissage auto-supervisé, méthodologie expérimentale, benchmark,
               robustesse, coût de calcul, Python, PyTorch, scikit-learn, statistiques, analyse de données,
               systèmes multi-agents, Reinforcement Learning, graphes, Graph Neural Networks, RAG,
               IA de confiance, explicabilité, cybersécurité, finance quantitative,
               Bac+5, Diplôme d'Ingénieur, école d'ingénieur, Grande École, stage de fin d'études},
  pdfborder={0 0 0}           % pas de cadre autour des liens (le soulignement bleu suffit)
}
% ==========================

\pagestyle{empty}

% Adjust margins
\addtolength{\oddsidemargin}{-0.45in}
\addtolength{\evensidemargin}{-0.45in}
\addtolength{\textwidth}{1.15in}
\addtolength{\topmargin}{-.62in}
\addtolength{\textheight}{1.20in}

\urlstyle{same}

\raggedbottom
\raggedright

%\definecolor{mydarkgreen}{RGB}{0, 140, 123}
\definecolor{mydarkgreen}{RGB}{0, 0, 0}

% Sections formatting
\titleformat{\section}{
  \vspace{-6pt}\scshape\raggedright\large
}{}{0em}{}[\color{mydarkgreen}\titlerule \vspace{-8pt}]

%-------------------------
% Custom commands

% Indented paragraph helper: \indentbox[ancho]{contenido}
% Uso: \indentbox{Texto...}  o  \indentbox[1.5em]{Texto...}
% [t] + \raggedright + \strut : lignes alignees en haut, sans espaces etires ni chevauchement.
\newcommand{\indentbox}[2][1em]{\hspace*{#1}\parbox[t]{\dimexpr\linewidth-#1\relax}{\raggedright #2\strut}}

% Photo et QR code desactives : dans Workday une photo casse l'analyse du texte voisin.


%-------------------------------------------
%%%%%%  CV STARTS HERE  %%%%%%%%%%%%%%%%%%%%%%%%%%%%


\begin{document}

%----------HEADING-----------------
% Nom seul sur la premiere ligne, puis coordonnees en texte brut dans le corps (pas d'en-tete, pas de minipage).

\noindent{\Large\bfseries\textcolor{mydarkgreen}{Javier Andres TARAZONA JIMENEZ}}\par
\vspace{3pt}
\noindent{\footnotesize javitar06@gmail.com \textbullet{} +33 7 46 36 97 75 \textbullet{} Palaiseau, France}\par
\vspace{1pt}
\noindent{\footnotesize\href{https://www.linkedin.com/in/javier-andres-tarazona-jimenez-84b489222/}{\textcolor{blue}{\underline{linkedin.com/in/javier-andres-tarazona-jimenez-84b489222}}} \textbullet{} \href{https://github.com/JavierTarazona06}{\textcolor{blue}{\underline{github.com/JavierTarazona06}}}}\par

\vspace{0.06cm}

{\centering
  {\large\bfseries AI Research Intern - Stage de fin d'études en Intelligence Artificielle}\\[2pt]
  {\footnotesize Diplôme d'Ingénieur Bac+5 en préparation (école d'ingénieur, Grande École, parcours IA) \textbullet{} Stage de 6 mois \textbullet{} Disponible à partir d'avril 2027}\par}

\vspace{0.05cm}

%-----------FORMATION-----------------
\section{\textcolor{mydarkgreen}{\textbf{FORMATION}}}
{\raggedright
\textbf{ENSTA Paris (Institut Polytechnique de Paris)} \textbullet{} Palaiseau, France \textbullet{} \textit{\footnotesize 2025/09 - 2027/09}\\
{\footnotesize\textit{Diplôme d'Ingénieur (Bac+5) - Master of Science in Engineering, parcours Intelligence Artificielle}}\\[4pt]
{\footnotesize Cours : Language Modeling (LLM), Deep Learning, Reinforcement Learning, Game Theory and Multi-Agent Control (systèmes multi-agents), Probability and Machine Learning, Graphes, jeux et recherche opérationnelle, Programmation GPU pour le Deep Learning, MLOps}\\[3pt]

\textbf{Universidad Nacional de Colombia (UNAL)} \textbullet{} Bogotá, Colombie \textbullet{} \textit{\footnotesize 2021/01 - 2025/07}\\
{\footnotesize\textit{Cursus d'ingénieur - Ingénierie des Systèmes et Informatique}}\\[4pt]
{\footnotesize 12e université d'Amérique latine (QS 2026). Cours : Probabilités et statistiques, Modèles stochastiques, Optimisation, Algèbre linéaire}\\[3pt]

\textbf{University of Saskatchewan} \textbullet{} Saskatoon, Canada \textbullet{} \textit{\footnotesize 2024/09 - 2024/12}\\
{\footnotesize\textit{Échange académique - Informatique : Algorithmes, Ingénierie logicielle, Systèmes d'exploitation}}\par}

\vspace{0.04cm}
%-----------EXPERIENCE-----------------
\section{\textcolor{mydarkgreen}{\textbf{EXPÉRIENCE PROFESSIONNELLE}}}
{\raggedright
\textbf{ENSTA Paris - Laboratoire U2IS} \textbullet{} Palaiseau, France \textbullet{} \textit{\footnotesize 2026/05 - 2026/08}\\
{\footnotesize\textit{Stagiaire de recherche en Intelligence Artificielle et Vision 3D}}\\[4pt]
\indentbox{\footnotesize \textbf{- Apprentissage auto-supervisé :} pré-entraînement de \textbf{Vision Transformers} (\textbf{PyTorch}) par prédiction de vues pour la structure 3D.}\\
\indentbox{\footnotesize \textbf{- Méthodologie expérimentale :} synthèse de nouvelles vues en curriculum learning sur formes 3D de complexité croissante.}\\
\indentbox{\footnotesize \textbf{- Pipeline de données en Python :} génération reproductible de jeux de données multi-vues avec Blender (rendu, caméras, métadonnées).}\\[3pt]

\textbf{APEDYS91} \textbullet{} Paris, France \textbullet{} \textit{\footnotesize 2025/10 - 2026/02}\\
{\footnotesize\textit{Ingénieur Machine Learning / NLP}}\\[4pt]
\indentbox{\footnotesize \textbf{- LLM et NLP :} correcteur de texte français en Python combinant règles (LanguageTool) et prompting de \textbf{LLM} locaux (Mistral, Gemma).}\\
\indentbox{\footnotesize \textbf{- Fiabilité et coût de calcul :} guardrails anti-hallucination (NER spaCy, entités, nombres) ; modèle choisi selon la RAM/VRAM.}\\[3pt]

\textbf{Engin A.I} \textbullet{} Bogotá, Colombie (à distance, États-Unis) \textbullet{} \textit{\footnotesize 2025/02 - 2025/06}\\
{\footnotesize\textit{Ingénieur Logiciel et Vision par Ordinateur}}\\[4pt]
\indentbox{\footnotesize \textbf{- Machine Learning et robustesse :} détection de voies (OpenCV, clustering non supervisé) : \textbf{+50\% de vitesse, +60\% de précision}.}\\
\indentbox{\footnotesize \textbf{- Analyse de données :} trajectoires de roue (WheelPath) corrélées aux zones de détérioration (\textbf{+10\% de satisfaction utilisateur}).}\\
\indentbox{\footnotesize \textbf{- Qualité du code Python :} refactorisation modulaire de 2 bibliothèques, \textbf{+70\% de score qualité} (maintenabilité, couverture de tests).}\\[3pt]

\textbf{Engin A.I} \textbullet{} Bogotá, Colombie (à distance, États-Unis) \textbullet{} \textit{\footnotesize 2024/01 - 2024/07}\\
{\footnotesize\textit{Ingénieur Logiciel et Vision par Ordinateur}}\par}


\vspace{0.04cm}
%-----------PROJETS-----------------
\section{\textcolor{mydarkgreen}{\textbf{PROJETS}}}
{\raggedright
\textbf{"Tameo" - Vision embarquée pour bateau autonome} \textbullet{} \textit{\footnotesize 2025/09 - 2027/06}\\
{\footnotesize\textit{Équipe de 5, membre (2025-2026) puis chef de projet (2026-2027) - Monaco Energy Boat Challenge, AI Class - Python, PyTorch, YOLO}}\\[4pt]
\indentbox{\footnotesize \textbf{- Benchmark de modèles :} fine-tuning, data augmentation et comparaison de détecteurs ; mAP et F1 portés d'environ 80\% à 90\%.}\\[3pt]

\textbf{Segmentation de peau par apprentissage statistique}\\
{\footnotesize\href{https://github.com/JavierTarazona06/Vision_LocalFeatures_BayesianClassification_KMeans}{\textcolor{blue}{\underline{github.com/JavierTarazona06/Vision\_LocalFeatures\_BayesianClassification\_KMeans}}}}\\
\indentbox{\footnotesize \textbf{- Étude comparative et rapport technique :} Gaussian Naive Bayes, QDA (supervisés) et K-Means sur RGB/HSV/YCrCb (scikit-learn).}\par}


\vspace{0.04cm}
%-----------COMPETENCES TECHNIQUES-----------------
\section{\textcolor{mydarkgreen}{\textbf{COMPÉTENCES TECHNIQUES}}}
{\footnotesize
\setlength{\parindent}{0pt}\setlength{\parskip}{2pt}
\hangindent=1.4em\hangafter=1 \textbf{Recherche en IA :} Machine Learning, Deep Learning, architectures Transformer, Large Language Models (LLM),\newline Natural Language Processing (NLP), prompting, fine-tuning, apprentissage supervisé et auto-supervisé, Computer Vision\par
\hangindent=1.4em\hangafter=1 \textbf{Méthodologie expérimentale :} benchmark et comparaison de modèles, métriques (précision, rappel, F1, mAP), robustesse, coût de calcul\par
\hangindent=1.4em\hangafter=1 \textbf{Statistiques et mathématiques :} probabilités et statistiques, analyse de données, modèles stochastiques, optimisation, algèbre linéaire\par
\hangindent=1.4em\hangafter=1 \textbf{Python et développement scientifique :} Python, PyTorch, scikit-learn, NumPy, pandas, Matplotlib, spaCy, OpenCV, Mistral, Gemma\par
\hangindent=1.4em\hangafter=1 \textbf{Ingénierie :} Git/GitHub, Docker, CI/CD, FastAPI, SQL, Google Cloud (GCP), C++/C, Java, Agile (SCRUM)\par
}

\vspace{0.04cm}
%-----------CENTRES D'INTERET-----------------
\section{\textcolor{mydarkgreen}{\textbf{CENTRES D'INTÉRÊT}}}
{\footnotesize
\noindent Graph Neural Networks (GNN) \textbullet{} Reinforcement Learning (apprentissage par renforcement) \textbullet{} IA de confiance et explicabilité \textbullet{} Retrieval-Augmented Generation (RAG) et recherche d'information \textbullet{} cybersécurité \textbullet{} finance quantitative et données financières\par
}

\vspace{0.04cm}
%-----------LANGUES-----------------
\section{\textcolor{mydarkgreen}{\textbf{LANGUES}}}
{\footnotesize
\noindent \textbf{Anglais :} courant (C1 - IELTS 2024) \textbullet{} \textbf{Français :} courant (B2 - DELF 2024) \textbullet{} \textbf{Espagnol :} langue maternelle \textbullet{} \textbf{Portugais :} notions\par
}

\vspace{0.04cm}
%--------DISTINCTIONS------------
\section{\textcolor{mydarkgreen}{\textbf{DISTINCTIONS}}}
{\footnotesize
\noindent \textbf{Bourse Eiffel Excellence} (ENSTA Paris) \textbullet{} \textbf{Mention d'Honneur - ICPC Colombie 2023} \textbullet{} \textbf{Excellence académique} (UNAL)\par
}


%-------------------------------------------
\end{document}
